\section{An explicit polylogarithm realization of the convolution calculus}
\label{poly:sec}

The circle Fourier description gives a direct bridge to the
polylogarithm's order derivatives at nonpositive integers. The underlying
Fourier summation is classical \cite{DLMFPolylog,DLMFHurwitz}.
The full mixed kernel and the symmetric first-Stieltjes specialization
also occur in the archived continuations
\cite{RepoConvolutionAlgebra,RepoConvolutionCalculus}.
This section consolidates their order-derivative form while preserving
the endpoint distributions.

\subsection{A two-parameter kernel}

For \(\alpha,\beta\in\bbC\), define the zero-mean distribution
\(\mathcal N_{\alpha,\beta}\) by
\begin{equation}\label{poly:N-fourier}
 \widehat{\mathcal N_{\alpha,\beta}}(n)=
 (2\pi\numi n)^\alpha(-2\pi\numi n)^\beta\quad(n\ne0),
 \qquad \widehat{\mathcal N_{\alpha,\beta}}(0)=0.
\end{equation}
The two powers use the logarithms specified in Section~\ref{scope:sec}.
On every compact set of parameters their coefficients and any fixed
number of parameter derivatives have a common polynomial growth
bound. This proves that \(\mathcal N\) is entire as a
distribution-valued function.

\begin{theorem}[Polylogarithm kernel for every mixed order]
\label{poly:master}
For \(0<h<1\), put \(z=\rme^{2\pi\numi h}\). Then
\begin{align}\label{poly:N}
 \mathcal N_{\alpha,\beta}(h)
 =(2\pi)^{\alpha+\beta}\bigl[
 &\rme^{\pi\numi(\alpha-\beta)/2}
                 \rLi_{-\alpha-\beta}(z)\notag\\
 &+\rme^{-\pi\numi(\alpha-\beta)/2}
                 \rLi_{-\alpha-\beta}(z^{-1})\bigr].
\end{align}
The formula represents the restriction to the open interval of the
entire distribution family, not an assertion of ordinary convergence
of its Fourier series at every order.

Let \(P_k\) be the spectral finite-part polygamma distribution, so that
\[
 \widehat P_k(n)=(2\pi\numi n)^k
  \bigl[\gamma+\log(2\pi\numi n)-H_k\bigr]\quad(n\ne0),
 \qquad H_0=0.
\]
For all \(k,l\geq0\), one has the distribution identity
\begin{equation}\label{poly:mixed}
 P_k*\check P_l=
 \left.
 (\partial_\alpha+\gamma-H_k)
 (\partial_\beta+\gamma-H_l)
       \mathcal N_{\alpha,\beta}
 \right|_{\alpha=k,\ \beta=l}.
\end{equation}
For the derivative-compatible normalization \(Q_k=D^kP_0\), the same
formula holds with \(H_k,H_l\) replaced by zero.
\end{theorem}

\begin{proof}
If \(\Re(\alpha+\beta)<-1\), the Fourier series is absolutely convergent.
The positive and negative indices separately give
\eqref{poly:N}. The same argument applies after parameter
differentiation on any smaller open region. The right side has an
entire continuation in its orders for \(z\ne1\), and is smooth in \(h\)
on compact subintervals of \((0,1)\). To identify its distributional
extension, insert the Abel factor \(\rho^{|n|}\), \(0<\rho<1\),
and take its limit in distributions. The Hurwitz Fourier formula
identifies the same local analytic continuation. The identity
therefore extends to all \(\alpha,\beta\). This argument also
specifies terms supported at \(h=0\), which a pointwise formula alone
could not recover.

Each parameter derivative in \eqref{poly:N-fourier} multiplies the
coefficient by the corresponding logarithm. At \((k,l)\), this gives
\(\widehat P_k(n)\widehat P_l(-n)\), exactly the Fourier coefficient of
\(P_k*\check P_l\). Fourier coefficients determine a distribution on
the circle. The formula for \(Q_k\) follows from its Fourier coefficient
\((2\pi\numi n)^k[\gamma+\log(2\pi\numi n)]\).
\end{proof}

The sine phases must be retained for complex parameters. A real-part
notation would not preserve holomorphic dependence on
\(\alpha,\beta\). At integer orders
\(\mathcal N_{k,l}=(-1)^lD^{k+l}\Delta\).
For \(k+l>0\) this is supported at the endpoint; for \(k=l=0\) its
open-interval value is \(-1\). Its parameter derivatives retain
additional, nonconstant interior terms. Substituting an integer
before differentiating would lose the desired information.

The unreflected kernel is obtained using one spectral parameter:
\(\mathcal M_\alpha=\mathcal N_{\alpha,0}\). Then
\begin{align}\label{poly:unreflected}
 P_k*P_l=\left.
 \bigl[\partial_\alpha^2
  +(2\gamma-H_k-H_l)\partial_\alpha
  +(\gamma-H_k)(\gamma-H_l)\bigr]\mathcal M_\alpha
 \right|_{\alpha=k+l}.
\end{align}
The equality is checked on each nonzero Fourier coefficient.
It also determines all the contact terms.

\subsection{First-Stieltjes identities at polylogarithm order zero}

Write \(c=\gamma+\log(2\pi)\), \(d=\pi/2\), and
\[
 L_j(h)=\left.\partial_s^j
                 \rLi_s(\rme^{2\pi\numi h})\right|_{s=0}.
\]
These derivatives are values of the analytically continued function.
They are not defined by directly summing a growing logarithmic
Fourier series.

\begin{corollary}[Two explicit order-derivative reductions]
\label{poly:gamma1}
For every \(0<h<1\),
\begin{align}
 \gamma_1(h)+\gamma_1(1-h)
 &=2\Re L_2(h)-4c\Re L_1(h)-c^2+\frac{\pi^2}{12},
 \label{poly:symmetric}\\
 \gamma_1(h)
 &=\Re L_2(h)-2c\Re L_1(h)+\pi\Im L_1(h)
      -\frac{c^2}{2}+\frac{\pi^2}{24}
      -\frac{\pi c}{2}\cot(\pi h).
 \label{poly:individual}
\end{align}
\end{corollary}

\begin{proof}
At \(k=l=0\), the operator on the positive-frequency part of
\eqref{poly:N} is
\[
 (c+d\numi-\partial_s)(c-d\numi-\partial_s)
       =(c-\partial_s)^2+d^2 .
\]
For real \(h\) the negative-frequency part is its conjugate.
The reflected convolution identity proved above reads on the open
interval
\[
 \gamma_1(h)+\gamma_1(1-h)-\frac{\pi^2}{3}
 =2\Re\bigl[L_2-2cL_1+(c^2+\pi^2/4)L_0\bigr].
\]
Since \(\Re L_0=-1/2\), this is \eqref{poly:symmetric}.

For the unreflected convolution, the operator is
\((c+d\numi-\partial_s)^2\). Its value is
\(2\gamma_1(h)+\zeta(2)\). Finally,
\[
 L_0(h)=-\frac12+\frac{\numi}{2}\cot(\pi h)
\]
gives \eqref{poly:individual} after expanding the real part.
Every constant comes from an explicit Gamma coefficient or
a contact term.
\end{proof}

These formulas belong to the classical differentiated Hurwitz
functional-equation framework. Their role here is an explicit
application of the finite-part normalization. They also check the
sign of the \(\cot(\pi h)\) contribution and the \(\pi^2\) constants.

At \(h=1/2\), the identity
\(\rLi_s(-1)=(2^{1-s}-1)\zeta(s)\) turns the corollary into an identity
among \(\zeta^{[j]}(0)\), \(\gamma_1\), \(\gamma\), and \(\log2\).
On the Stieltjes side,
\[
 \gamma_1(1/2)=\gamma_1-2\gamma\log2-\log^22,
\]
obtained directly by expanding
\(\zeta(s,1/2)=(2^s-1)\zeta(s)\).
Thus the convergent digamma convolution from the preceding part,
at \(a=1/2\), has the explicit value
\begin{equation}\label{poly:half-convolution}
 2\gamma_1-4\gamma\log2-2\log^22+\frac{\pi^2}{6}.
\end{equation}
This evaluates the fully subtracted integral; the corresponding
unsubtracted product integral diverges.

\subsection{Rational shifts and a finite coordinate algorithm}

The pending report asks for rational-shift coordinate formulas
\cite{RepoIncoming}. The conversion follows from a finite Fourier
transform already present in the manuscript. The distributional
extension now makes it applicable to mixed polygamma kernels at
singular spectral orders.

For \(q\geq2\), \(1\leq p<q\), and \(z=\rme^{2\pi\numi p/q}\),
\begin{equation}\label{poly:rational}
 \rLi_s(z)=q^{-s}\sum_{a=1}^q z^a\zeta(s,a/q).
\end{equation}
At any regular spectral point \(s_0\),
\begin{equation}\label{poly:rational-jets}
 \rLi_{s_0}^{[j]}(z)
 =q^{-s_0}\sum_{a=1}^q z^a
   \sum_{v=0}^j\binom jv(-\log q)^{j-v}
             \zeta^{[v]}(s_0,a/q).
\end{equation}
At \(s=1\) the pole cancels because \(\sum_{a=1}^qz^a=0\);
one must expand the finite combined expression before specializing.
Equations \eqref{poly:N}--\eqref{poly:rational-jets}
are a finite algorithm for every rational-shift polygamma
convolution and every fixed spectral derivative.

A multiplicative-character transform then converts these Hurwitz
coordinates to primitive Dirichlet-\(L\) coordinates. Conductor
logarithms are generated by differentiating the explicit factors.
The manuscript's conductor-lifted trace theorem and pole-corrected
jets already prove that arithmetic conversion \cite{RepoManuscript}.
We use them as an available normal form and do not claim
independence of the resulting evaluated constants.
